% math0 manuscript. Build with: latexmk -pdf paper.tex
\documentclass[11pt]{amsart}
\usepackage[margin=1.1in]{geometry}
\usepackage{amsmath,amssymb,amsthm,mathtools}
\usepackage{url}
\usepackage{hyperref}
\usepackage{xcolor}
\usepackage{listings}
\usepackage{newunicodechar}

\lstdefinelanguage{lean}{
  keywords={theorem,lemma,def,abbrev,example,fun,by,have,show,from,let,in,if,then,else,match,with,
    namespace,end,open,import,section,variable,instance,structure,where,calc,obtain,intro,exact,
    refine,rcases,rintro,simp,omega,decide,norm_num,linarith,nlinarith,positivity,induction,cases,
    private,noncomputable},
  sensitive=true, comment=[l]{--}, morecomment=[s]{/-}{-/}, string=[b]"}
\lstset{language=lean, basicstyle=\ttfamily\footnotesize, columns=fullflexible, keepspaces=true,
  breaklines=true, keywordstyle=\color{blue!60!black}\bfseries, commentstyle=\color{gray},
  frame=single, extendedchars=true, inputencoding=utf8,
  literate={∑}{{$\sum$}}1 {∈}{{$\in$}}1 {∉}{{$\notin$}}1 {ℕ}{{$\mathbb{N}$}}1 {ℤ}{{$\mathbb{Z}$}}1
    {∀}{{$\forall$}}1 {∃}{{$\exists$}}1 {∧}{{$\wedge$}}1 {≤}{{$\le$}}1 {→}{{$\to$}}1
    {×}{{$\times$}}1 {²}{{$^2$}}1 {³}{{$^3$}}1}
% Unicode commonly found in Lean sources.
\newunicodechar{ℕ}{\ensuremath{\mathbb{N}}}
\newunicodechar{ℤ}{\ensuremath{\mathbb{Z}}}
\newunicodechar{ℚ}{\ensuremath{\mathbb{Q}}}
\newunicodechar{ℝ}{\ensuremath{\mathbb{R}}}
\newunicodechar{ℂ}{\ensuremath{\mathbb{C}}}
\newunicodechar{∀}{\ensuremath{\forall}}
\newunicodechar{∃}{\ensuremath{\exists}}
\newunicodechar{→}{\ensuremath{\to}}
\newunicodechar{←}{\ensuremath{\leftarrow}}
\newunicodechar{↔}{\ensuremath{\leftrightarrow}}
\newunicodechar{↦}{\ensuremath{\mapsto}}
\newunicodechar{¬}{\ensuremath{\neg}}
\newunicodechar{∧}{\ensuremath{\wedge}}
\newunicodechar{∨}{\ensuremath{\vee}}
\newunicodechar{≤}{\ensuremath{\le}}
\newunicodechar{≥}{\ensuremath{\ge}}
\newunicodechar{≠}{\ensuremath{\ne}}
\newunicodechar{∈}{\ensuremath{\in}}
\newunicodechar{∉}{\ensuremath{\notin}}
\newunicodechar{⊆}{\ensuremath{\subseteq}}
\newunicodechar{∩}{\ensuremath{\cap}}
\newunicodechar{∪}{\ensuremath{\cup}}
\newunicodechar{∑}{\ensuremath{\sum}}
\newunicodechar{∏}{\ensuremath{\prod}}
\newunicodechar{⟨}{\ensuremath{\langle}}
\newunicodechar{⟩}{\ensuremath{\rangle}}
\newunicodechar{·}{\ensuremath{\cdot}}
\newunicodechar{∘}{\ensuremath{\circ}}
\newunicodechar{⁻}{\ensuremath{^{-}}}
\newunicodechar{¹}{\ensuremath{^{1}}}
\newunicodechar{²}{\ensuremath{^{2}}}
\newunicodechar{₀}{\ensuremath{_{0}}}
\newunicodechar{₁}{\ensuremath{_{1}}}
\newunicodechar{₂}{\ensuremath{_{2}}}
\newunicodechar{α}{\ensuremath{\alpha}}
\newunicodechar{β}{\ensuremath{\beta}}
\newunicodechar{ε}{\ensuremath{\varepsilon}}
\newunicodechar{δ}{\ensuremath{\delta}}
\newunicodechar{π}{\ensuremath{\pi}}
\newunicodechar{σ}{\ensuremath{\sigma}}
\newunicodechar{φ}{\ensuremath{\varphi}}
\newunicodechar{∞}{\ensuremath{\infty}}
\newunicodechar{⌊}{\ensuremath{\lfloor}}
\newunicodechar{⌋}{\ensuremath{\rfloor}}
\newunicodechar{⌈}{\ensuremath{\lceil}}
\newunicodechar{⌉}{\ensuremath{\rceil}}
\newunicodechar{∣}{\ensuremath{\mid}}
\newunicodechar{≡}{\ensuremath{\equiv}}
\newunicodechar{×}{\ensuremath{\times}}
\newunicodechar{⧸}{\ensuremath{/}}
\newunicodechar{‖}{\ensuremath{\|}}
\newunicodechar{√}{\ensuremath{\surd}}
\newunicodechar{∅}{\ensuremath{\emptyset}}
\newunicodechar{⦃}{\ensuremath{\{\!|}}
\newunicodechar{⦄}{\ensuremath{|\!\}}}
\newunicodechar{𝓝}{\ensuremath{\mathcal{N}}}
\newunicodechar{≫}{\ensuremath{\gg}}
\newunicodechar{≪}{\ensuremath{\ll}}

\theoremstyle{plain}
\newtheorem{theorem}{Theorem}
\newtheorem{lemma}[theorem]{Lemma}
\newtheorem{proposition}[theorem]{Proposition}
\newtheorem{corollary}[theorem]{Corollary}
\newtheorem{conjecture}[theorem]{Conjecture}
\theoremstyle{definition}
\newtheorem{definition}[theorem]{Definition}
\newtheorem{remark}[theorem]{Remark}

\emergencystretch=3em
% Lean identifiers in running text (breakable at dots and underscores).
\expandafter\def\expandafter\UrlBreaks\expandafter{\UrlBreaks\do\_\do\.}
\DeclareUrlCommand\decl{\urlstyle{tt}}
\newcommand{\N}{\mathbb{N}}
\newcommand{\Z}{\mathbb{Z}}
\newcommand{\D}{\mathcal{D}}
\newcommand{\LB}{\mathrm{LB}}
\newcommand{\ind}[1]{[\,#1\,]}

\title[A polycube of $k^3$ dice shows at least $6k(k+2)$ pips]{For $k\ge2$, a face-connected polycube of $k^3$ standard dice shows at least $6k(k+2)$ pips, with equality for the cube (OEIS A363609)}
\author{math0}
\address{Produced by Claude (an AI model developed by Anthropic) in the math0 project}
\date{October 19, 2026}

\begin{document}
\begin{abstract}
Let $a(n)$ be the least total number of pips that can be visible on a face-connected polycube built from $n$ standard dice, where a face of a die is visible unless it is glued to another die (OEIS A363609). We prove the conjecture recorded in that entry: $a(k^3)=6k(k+2)$ for every $k\ge2$, so the $k\times k\times k$ cube is optimal among all polycubes with $k^3$ cells. The lower bound combines a per-die inequality (a die shows at least as many pips as it has exposed faces plus exposed pairs of perpendicular faces) with projection counts and a slice-wise discrete Loomis--Whitney inequality; it does not use connectivity. The math0 referees found no earlier resolution, and we classify the result as new. The proof is verified in Lean~4 against a statement that math0 formalized from the OEIS entry, reviewed for fidelity and frozen before the proof attempt; the Lean kernel checks its definitions against the entry's data, exactly for the first seven terms and as upper bounds for all $57$ listed terms.
\end{abstract}
\maketitle

\section{Introduction}

A standard die carries $1,\dots,6$ pips on its six faces, with opposite faces summing to $7$. Glue $n$ such dice face to face into a face-connected polycube and count the pips on all faces that are not glued to another die, including the bottom faces. The sequence A363609 \cite{A363609}, contributed by Matt Donahoe on June~11, 2023, lists the least possible count $a(n)$ for $n=1,\dots,57$:
\[
21,30,40,40,51,52,54,48,60,62,65,60,\dots,90\ (n=27),\dots
\]
The entry's formula section contains the line \texttt{Conjecture: a(k\^{}3) = 6*(k+2)*k for k > 1.} For the cube $[0,k)^3$ the value $6k(k+2)$ is attained: turn the $8$ corner dice to show $1,2,3$, the $12(k-2)$ edge dice to show $1,2$ and the $6(k-2)^2$ face dice to show $1$, for a total of $48+36(k-2)+6(k-2)^2=6k(k+2)$. The content of the conjecture is the matching lower bound over all polycubes with $k^3$ cells.

\subsection*{Prior work}
The entry links a Python program by Donahoe \cite{Don}, which the referees describe as a heuristic search over layouts consisting of a box, a rectangular partial layer and a row; it does not prove optimality. M.~S.~Branicky added two formulas to the entry in June 2023: the upper bound $a(ijk)\le 48+2((i-2)(j-2)+(i-2)(k-2)+(j-2)(k-2))+12(i+j+k-6)$ for $i,j,k>1$, which for $i=j=k$ is the cube count above, and two-sided bounds $\mathrm{lb}(n)\le a(n)\le\mathrm{ub}(n)$ in terms of the minimum surface area A193416 \cite{A193416}. These do not settle the conjecture. The value $90$ for the $3\times3\times3$ cube appears as a contest exercise (AMC~10A 2002, Problem~18 \cite{AMC}); it concerns the cube alone, not optimality among polycubes. The unweighted analogue is A193416, the minimum surface area of a polycube with $n$ cells, for which the entry records $a(n^3)=6n^2$ (C.~R.~Greathouse~IV, 2011); the lower bound $6n^2$ follows from the Loomis--Whitney inequality \cite{LW} as in Section~\ref{sec:proj}. The math0 referees found the conjecture still stated as open in the current revision (\#45, September~3, 2023) of A363609 and found no proof or counterexample in searches of arXiv, MathOverflow, Mathematics Stack Exchange, Puzzling Stack Exchange and SeqFan. A discussion that is not indexed cannot be excluded, but on this evidence the result is new.

Besides the classical projection bound for the surface area, our proof of the conjecture (Theorem~\ref{thm:main}) uses an \emph{edge term}: a die shows at least one extra pip for each pair of perpendicular exposed faces, and every slice of the polycube has a cell with two prescribed perpendicular faces exposed. The proof is formalized in Lean~4 with Mathlib (Section~\ref{sec:lean}).

\section{Statement}\label{sec:statement}

Write $e_1,e_2,e_3$ for the standard basis of $\Z^3$ and $\D=\{\pm e_1,\pm e_2,\pm e_3\}$. A cell $q\in\Z^3$ is the unit cube $q+[0,1]^3$; its six faces have outer normals $d\in\D$. Two directions $d,e\in\D$ are \emph{perpendicular} if $e\ne\pm d$.

\begin{definition}\label{def:main}
A \emph{polycube} is a finite set $P\subset\Z^3$ that is face-connected: any two cells of $P$ are joined by a chain of cells of $P$ in which consecutive cells differ by an element of $\D$. An \emph{orientation} is a pair $o=(u,v)$ of perpendicular $u,v\in\D$; a die in orientation $o$ shows $\pi_o(d)$ pips on its face with normal $d$, where
\[
\pi_o(u)=1,\quad \pi_o(v)=2,\quad \pi_o(u\times v)=3,\quad \pi_o(-u\times v)=4,\quad \pi_o(-v)=5,\quad \pi_o(-u)=6.
\]
For a polycube $P$ and orientations $o=(o_q)_{q\in P}$, the number of visible pips is
\[
V(P,o)=\sum_{q\in P}\ \sum_{d\in\D,\ q+d\notin P}\pi_{o_q}(d),
\]
and $a(n)$ is the minimum of $V(P,o)$ over all polycubes $P$ with $|P|=n$ and all $o$.
\end{definition}

There are $6\cdot4=24$ orientations, the rotations of one die; opposite faces sum to $7$, and $1,2,3$ form a right-handed frame, so all dice have the same chirality. A reflection of a whole configuration replaces every die by its mirror image, so the other chirality gives the same minimum.

\begin{theorem}\label{thm:main}
For every integer $k\ge2$, $a(k^3)=6k(k+2)$.
\end{theorem}

\subsection*{The OEIS entry}
The entry A363609 \cite{A363609} (author Matt Donahoe, June~11, 2023; revision \#45, September~3, 2023) has the name ``Minimum sum of the visible pips on a polycube made from n dice.'', offset $1$, the comment ``This sequence is calculated using standard six-sided dice of the same chirality. Opposite sides sum to seven.'' and the example ``For n = 2, two dice are conjoined to hide both their 6-pip faces, so a(2) = 30.'' Its formula section contains, verbatim and without attribution, the line
\begin{quote}\raggedright
\texttt{Conjecture: a(k\^{}3) = 6*(k+2)*k for k > 1.}
\end{quote}
Theorem~\ref{thm:main} is this conjecture.

\subsection*{Lean formalization}
The statement was formalized by math0 in \path{lean/Math0/Statements/OEIS/A363609.lean}, reviewed for fidelity and frozen (SHA-256 \texttt{aec2c052\ldots 245afe}) before any proof attempt; it is not part of the formal-conjectures repository \cite{FC}. Its definitions and the statement are, verbatim except that docstrings and the certification section are omitted:
\begin{lstlisting}
namespace Math0.Statements.A363609
open Finset
abbrev Cell := ℤ × ℤ × ℤ
def dirs : Finset Cell := {(1, 0, 0), (-1, 0, 0), (0, 1, 0), (0, -1, 0), (0, 0, 1), (0, 0, -1)}
def dot (u v : Cell) : ℤ := u.1 * v.1 + u.2.1 * v.2.1 + u.2.2 * v.2.2
def cross (u v : Cell) : Cell :=
  (u.2.1 * v.2.2 - u.2.2 * v.2.1, u.2.2 * v.1 - u.1 * v.2.2, u.1 * v.2.1 - u.2.1 * v.1)
def IsOrientation (o : Cell × Cell) : Prop := o.1 ∈ dirs ∧ o.2 ∈ dirs ∧ dot o.1 o.2 = 0
def pips (o : Cell × Cell) (d : Cell) : ℕ :=
  if d = o.1 then 1 else if d = o.2 then 2 else if d = cross o.1 o.2 then 3
  else if d = -cross o.1 o.2 then 4 else if d = -o.2 then 5 else if d = -o.1 then 6 else 0
def visiblePips (P : Finset Cell) (o : Cell → Cell × Cell) : ℕ :=
  ∑ c ∈ P, ∑ d ∈ dirs, if c + d ∈ P then 0 else pips (o c) d
def Adj (P : Finset Cell) (u v : Cell) : Prop := u ∈ P ∧ v ∈ P ∧ v - u ∈ dirs
def FaceConnected (P : Finset Cell) : Prop :=
  ∀ x ∈ P, ∀ y ∈ P, Relation.ReflTransGen (Adj P) x y
def totals (n : ℕ) : Set ℕ :=
  {t | ∃ P : Finset Cell, P.card = n ∧ FaceConnected P ∧
    ∃ o : Cell → Cell × Cell, (∀ c ∈ P, IsOrientation (o c)) ∧ visiblePips P o = t}
noncomputable def a (n : ℕ) : ℕ := sInf (totals n)
def statement : Prop := ∀ k : ℕ, 1 < k → a (k ^ 3) = 6 * (k + 2) * k
end Math0.Statements.A363609
\end{lstlisting}
The encoding matches Definition~\ref{def:main}. \decl{dirs} is $\D$, \decl{IsOrientation} says that $u,v\in\D$ are perpendicular, and \decl{pips} is $\pi_o$. The face of $c\in P$ with normal $d$ counts iff $c+d\notin P$, so all six faces of a single die count and $a(1)=21$. A face towards an enclosed empty cell also counts; this agrees with the data, since under the other reading the $3\times3\times3$ cube without its centre would give $a(26)\le90$, whereas the entry lists $a(26)=93$. (See also Remark~\ref{rem:extra}.) \texttt{totals n} is the set of all values $V(P,o)$ and \texttt{a n} its infimum in $\N$. For $n\ge1$ the set is nonempty (a row of $n$ cells is a polycube), so the infimum is a minimum; the value \texttt{a 0} is not part of the sequence and is excluded by $k>1$. The statement involves no truncated subtraction or division.

The statement file checks the definitions against the data by kernel computation (\texttt{decide +kernel}). The theorem \decl{a_data} gives the exact values $a(1),\dots,a(7)=21,30,40,40,51,52,54$, by a certified enumeration of all polycubes with at most $6$ cells up to the $48$ lattice symmetries ($1,1,2,7,23,112$ representatives) and of all their one-cell extensions. The theorem \decl{a_le_data} gives $a(n)\le$ the listed value for all $57$ listed terms, via kernel-checked witness polycubes. Unformalized exhaustive enumerations run during review reproduced the numbers of fixed polycubes (A001931 \cite{A001931}) and the listed values up to $a(9)$ or $a(10)$.

\section{Proof}\label{sec:proof}

The proof follows \path{lean/Math0/Frontier/OEIS/A363609.lean}; Section~\ref{sec:lean} lists the corresponding declarations. Throughout, $P\subset\Z^3$ is finite and $q=(q_1,q_2,q_3)$ denotes a cell.

\subsection{Reduction to one die at a time}
For $q\in P$ let $E_P(q)=\{d\in\D: q+d\notin P\}$ be the set of exposed directions. For $E\subseteq\D$ let $f(E)=\min_o\sum_{d\in E}\pi_o(d)$, the minimum over the $24$ orientations, and put
\[
\LB(P)=\sum_{q\in P}f(E_P(q)).
\]

\begin{lemma}\label{lem:reduce}
For every $n$: if $P_0$ is a polycube with $|P_0|=n$ and $\LB(P_0)\le\LB(P)$ for every polycube $P$ with $|P|=n$, then $a(n)=\LB(P_0)$.
\end{lemma}

\begin{proof}
For fixed $P$ the term of $V(P,o)$ for the cell $q$ depends only on $o_q$, so $V(P,o)\ge\LB(P)$ for every $o$, with equality when each $o_q$ is chosen to minimize $\sum_{d\in E_P(q)}\pi_{o_q}(d)$. Hence $\LB(P_0)$ is a value of $V$ and is at most every value of $V$ on polycubes with $n$ cells.
\end{proof}

For $E\subseteq\D$ let $p(E)$ be the number of unordered pairs $\{d,e\}\subseteq E$ of perpendicular directions.

\begin{lemma}\label{lem:die}
For every $E\subseteq\D$, $f(E)\ge|E|+p(E)$, with equality if $E$ contains no two opposite directions.
\end{lemma}

\begin{proof}
Fix an orientation. The faces in $E$ show distinct numbers of pips, and two opposite faces show $7$ together. Let $r$ be the number of opposite pairs $\{d,-d\}\subseteq E$ and $s=|E|-2r$ the number of remaining faces of $E$; these lie in $s$ further opposite pairs, so $r+s\le3$, and they show $s$ distinct positive numbers. Hence the pips on $E$ total at least $7r+s(s+1)/2$, while $|E|+p(E)=(2r+s)+\binom{2r+s}{2}-r$. For $r=0$ both sides equal $s(s+1)/2$. The remaining cases are $(r,s)=(1,0),(1,1),(1,2),(2,0),(2,1),(3,0)$, where the two sides are $7\ge2$, $8\ge5$, $10\ge9$, $14\ge8$, $15\ge13$ and $21\ge18$.

If $r=0$, the $s\le3$ faces of $E$ are mutually perpendicular. Choosing $u,v$ among them (in the right order, so that $u\times v$ is the third one when $s=3$) puts $1,\dots,s$ on them, so $f(E)\le s(s+1)/2=|E|+p(E)$.
\end{proof}

In Lean, $f$ is evaluated by the kernel on all $64$ subsets $E$ (minimum over the $24$ orientations), giving a table \decl{tabF}; the inequality and the equality case are then checked on that table by \texttt{decide}. For $d\in\D$ and perpendicular $d,e\in\D$ let
\[
N_d(P)=\#\{q\in P: q+d\notin P\},\qquad N_{d,e}(P)=\#\{q\in P: q+d\notin P,\ q+e\notin P\}.
\]
Summing Lemma~\ref{lem:die} over $q\in P$ gives
\begin{equation}\label{eq:LBE}
\LB(P)\ \ge\ \sum_{d\in\D}N_d(P)+\sum_{\{d,e\}}N_{d,e}(P),
\end{equation}
the second sum running over the $12$ unordered pairs of perpendicular directions.

\subsection{Projections and slices}\label{sec:proj}
For $i\in\{1,2,3\}$ let $A_i$ be the number of points of the projection of $P$ that forgets the coordinate $q_i$, and let $\ell_i=\#\{q_i:q\in P\}$ be the number of occupied slices orthogonal to $e_i$.

\begin{lemma}\label{lem:faces}
$N_{\pm e_i}(P)\ge A_i$ for every $i$.
\end{lemma}

\begin{proof}
Let $d=se_i$ with $s=\pm1$, and let $\rho$ be the projection forgetting $q_i$. For each point $w\in\rho(P)$ choose $m\in P\cap\rho^{-1}(w)$ with $sm_i$ maximal. Then $m+d$ lies in the same fibre and has a larger value of $sq_i$, so $m+d\notin P$. Thus $\rho$ maps $\{q\in P:q+d\notin P\}$ onto $\rho(P)$.
\end{proof}

\begin{lemma}\label{lem:edges}
Let $\{i,j,l\}=\{1,2,3\}$ and $s,t\in\{\pm1\}$. Then $N_{se_i,te_j}(P)\ge\ell_l$.
\end{lemma}

\begin{proof}
For each occupied value $z$ of $q_l$ choose $m\in P$ with $m_l=z$ such that $(sm_i,tm_j)$ is lexicographically maximal. The cells $m+se_i$ and $m+te_j$ have $l$-th coordinate $z$; the first has a larger value of $sq_i$, the second the same value of $sq_i$ and a larger value of $tq_j$. Both exceed $(sm_i,tm_j)$ lexicographically, so neither lies in $P$. Thus $q\mapsto q_l$ maps $\{q\in P:q+se_i\notin P,\ q+te_j\notin P\}$ onto the set of occupied values.
\end{proof}

Each axis carries two directions and each of the three pairs of axes carries four perpendicular pairs of directions, so \eqref{eq:LBE} and Lemmas~\ref{lem:faces} and~\ref{lem:edges} give
\begin{equation}\label{eq:proj}
\LB(P)\ \ge\ 2(A_1+A_2+A_3)+4(\ell_1+\ell_2+\ell_3).
\end{equation}

\begin{lemma}\label{lem:slices}
If $|P|=k^3$, then $\ell_1+\ell_2+\ell_3\ge3k$.
\end{lemma}

\begin{proof}
$P$ is contained in the product of its three coordinate sets, so $k^3\le\ell_1\ell_2\ell_3$. By the AM--GM inequality $27\ell_1\ell_2\ell_3\le(\ell_1+\ell_2+\ell_3)^3$, hence $(3k)^3\le(\ell_1+\ell_2+\ell_3)^3$.
\end{proof}

\begin{lemma}[slice-wise Loomis--Whitney]\label{lem:LW}
For every $k\in\N$, $4k|P|\le(A_3+k^2)(A_1+A_2)$. In particular, if $|P|=k^3$ then $A_1+A_2+A_3\ge3k^2$.
\end{lemma}

\begin{proof}
For an occupied value $z$ of $q_3$ let $P_z=\{q\in P:q_3=z\}$, $m_z=|P_z|$, and let $\alpha_z,\beta_z$ be the sizes of the projections of $P_z$ forgetting $q_1$ and $q_2$, respectively. These projections keep $q_3$, so the projections of $P$ split by the value of $q_3$ and $\sum_z\alpha_z=A_1$, $\sum_z\beta_z=A_2$, $\sum_z m_z=|P|$. In $P_z$ a cell is determined by $(q_1,q_2)$, so $m_z\le A_3$; it is also determined by its two projections, so $m_z\le\alpha_z\beta_z$. Hence, using $4A_3k^2\le(A_3+k^2)^2$ and $4\alpha_z\beta_z\le(\alpha_z+\beta_z)^2$,
\[
(4km_z)^2=4k^2\cdot4m_z^2\le(4A_3k^2)(4\alpha_z\beta_z)\le\bigl((A_3+k^2)(\alpha_z+\beta_z)\bigr)^2,
\]
so $4km_z\le(A_3+k^2)(\alpha_z+\beta_z)$. Summing over $z$ gives the first claim. If $|P|=k^3$, put $u=A_3+k^2$ and $v=A_1+A_2$; then $16k^4\le4uv\le(u+v)^2$, so $u+v\ge4k^2$.
\end{proof}

The classical Loomis--Whitney inequality \cite{LW}, applied to the union of the unit cubes of $P$, gives the discrete form $|P|^2\le A_1A_2A_3$; with AM--GM it gives the same consequence. The form above, in which $\sqrt{A_3}$ is replaced by the AM--GM bound $(A_3+k^2)/(2k)$, avoids square roots.

\begin{proposition}\label{prop:lower}
If $P\subset\Z^3$ is finite with $|P|=k^3$, then $\LB(P)\ge6k(k+2)$.
\end{proposition}

\begin{proof}
By \eqref{eq:proj} and Lemmas~\ref{lem:slices} and~\ref{lem:LW}, $\LB(P)\ge2\cdot3k^2+4\cdot3k=6k^2+12k$.
\end{proof}

\subsection{The cube}

\begin{proposition}\label{prop:cube}
Let $k\ge2$ and $Q_k=\{0,\dots,k-1\}^3$. Then $Q_k$ is a polycube with $k^3$ cells and $\LB(Q_k)=6k(k+2)$.
\end{proposition}

\begin{proof}
Every cell of $Q_k$ is joined to the origin inside $Q_k$ by moving along the first, then the second, then the third axis, so $Q_k$ is face-connected. For $q\in Q_k$ the direction $e_i$ is exposed iff $q_i=k-1$, and $-e_i$ iff $q_i=0$; as $k\ge2$, no two opposite directions are exposed. Put $u(x)=\ind{x=k-1}+\ind{x=0}\in\{0,1\}$ and $X_i=u(q_i)$. By the equality case of Lemma~\ref{lem:die}, $f(E_{Q_k}(q))=X_1+X_2+X_3+X_1X_2+X_1X_3+X_2X_3$, that is,
\[
f(E_{Q_k}(q))+1+X_1X_2X_3=(1+X_1)(1+X_2)(1+X_3).
\]
A product of functions of one coordinate each sums over $Q_k$ to the product of the one-dimensional sums, and $\sum_{x=0}^{k-1}u(x)=2$. Summing over $q\in Q_k$ therefore gives $\LB(Q_k)+k^3+2^3=(k+2)^3$, so $\LB(Q_k)=6k^2+12k$.
\end{proof}

\begin{proof}[Proof of Theorem~\ref{thm:main}]
Let $k\ge2$. By Proposition~\ref{prop:cube}, $Q_k$ is a polycube with $k^3$ cells and $\LB(Q_k)=6k(k+2)$, and by Proposition~\ref{prop:lower} every polycube $P$ with $k^3$ cells has $\LB(P)\ge6k(k+2)$. Lemma~\ref{lem:reduce} gives $a(k^3)=6k(k+2)$.
\end{proof}

\begin{remark}\label{rem:extra}
The following observations are not part of the formal development.
(1) Connectivity of $P$ is used only for the upper bound.
(2) Lemma~\ref{lem:die} holds for any labelling of a die with opposite faces summing to $7$, so the argument does not depend on chirality.
(3) Every face counted in Lemmas~\ref{lem:faces} and~\ref{lem:edges} sees infinity along a straight ray that avoids $P$, so these lemmas and \eqref{eq:LBE} remain valid when only faces facing the unbounded component of the complement are counted. Hence Theorem~\ref{thm:main} also holds if faces towards an enclosed cavity are not counted.
(4) For general $n$, the inequality $|P|^2\le A_1A_2A_3$ \cite{LW}, $|P|\le\ell_1\ell_2\ell_3$, AM--GM and \eqref{eq:proj} give $a(n)\ge6n^{2/3}+12n^{1/3}$, with equality for $n=k^3$, $k\ge2$.
\end{remark}

\section{Machine verification}\label{sec:lean}

\subsection*{Files and declarations}
The proof is in \path{lean/Math0/Frontier/OEIS/A363609.lean} in the math0 repository, \url{https://github.com/Foundation-0/math0}. It imports the frozen statement file \path{lean/Math0/Statements/OEIS/A363609.lean} (Section~\ref{sec:statement}) and the audit commands in \path{lean/Math0/Audit.lean}. The main declaration is
\begin{lstlisting}
theorem result : statement
\end{lstlisting}
in the namespace \decl{Math0.Frontier.A363609}, so \decl{Math0.Frontier.A363609.result} has type exactly \decl{Math0.Statements.A363609.statement}. From the statement file the proof uses \decl{fmin_exposed} (the value $f(E_P(q))$ is the table entry \decl{tabF} at the six exposure flags, proved by kernel evaluation of the minimum over the $24$ orientations), \decl{LB_mem_totals} and \decl{a_eq_of} (together Lemma~\ref{lem:reduce}). The private auxiliary declarations of the proof file correspond to Section~\ref{sec:proof} as follows: \decl{tabF_lower} and \decl{tabF_upper} are Lemma~\ref{lem:die}, as $64$-case checks of the table by \texttt{decide}; \decl{E1}, \decl{E2} and \decl{LB_ge_E} are $N_d$, $N_{d,e}$ and \eqref{eq:LBE}; \decl{card_image_le_E1} and \decl{card_image_le_E2} are Lemmas~\ref{lem:faces} and~\ref{lem:edges} for an abstract projection and height function (via \decl{Finset.exists_max_image}, with the key \decl{toLex} in the second); \decl{LB_ge_proj} is \eqref{eq:proj}; \decl{amgm3} and \decl{sum3_ge} are Lemma~\ref{lem:slices}; \decl{slice_arith}, \decl{loomis} and \decl{sum2_ge} are Lemma~\ref{lem:LW}; \decl{LB_lower} is Proposition~\ref{prop:lower}; and \decl{cube}, \decl{card_cube}, \decl{mem_cube}, \decl{line}, \decl{cube_conn}, \decl{sum_cube_mul} and \decl{LB_cube} give Proposition~\ref{prop:cube}. Apart from small kernel-checked facts about the six directions, the only kernel computations the proof relies on are these two checks and, in the statement file, the evaluation of \decl{tabF} (\decl{fmin_sel}, by \texttt{decide +kernel}). \texttt{sorry}, \texttt{native\_decide} and \texttt{set\_option} do not occur in either file.

\subsection*{Audit}
The file ends with the command
\begin{lstlisting}
#math0_settles Math0.Frontier.A363609.result Math0.Statements.A363609.statement
\end{lstlisting}
from \path{lean/Math0/Audit.lean}. It fails unless (1) \decl{statement} is a definition of type \texttt{Prop} without universe parameters; (2) \decl{statement} and \decl{result} depend on no axioms other than \decl{propext}, \decl{Classical.choice} and \decl{Quot.sound} (this excludes \decl{sorryAx}, \texttt{native\_decide} and custom axioms); and (3) the type of \decl{result} is syntactically the statement or its negation, after removing metadata annotations, with no unfolding. On this file it reports
\begin{lstlisting}[basicstyle=\ttfamily\scriptsize]
math0 audit: 'Math0.Statements.A363609.statement' depends on axioms [Classical.choice, Quot.sound, propext]
math0 audit: 'Math0.Frontier.A363609.result' depends on axioms [Classical.choice, Quot.sound, propext]
math0 statement check: 'Math0.Frontier.A363609.result' proves 'Math0.Statements.A363609.statement'
\end{lstlisting}

\subsection*{Toolchain and reproduction}
We use Lean~4 \cite{Lean4} \texttt{v4.33.1} and Mathlib \cite{Mathlib} \texttt{v4.33.1}, pinned in \path{lean/}; the repository also pins formal-conjectures \cite{FC} at commit \texttt{c77fdd4}, which this result does not use. To reproduce, run \path{scripts/verify.sh} from the repository root; it fetches the Mathlib cache, builds all math0 results with \texttt{lake build} (re-running every audit command) and checks that the pinned dependencies are unmodified. The single file is checked by \texttt{lake env lean Math0/Frontier/OEIS/A363609.lean} in \path{lean}, in about $16$ seconds once the statement module (whose kernel data checks take a few minutes) is built.

\begin{thebibliography}{99}

\bibitem{AMC}
Art of Problem Solving, \emph{2002 AMC 10A Problems/Problem 18}, \url{https://artofproblemsolving.com/wiki/index.php/2002_AMC_10A_Problems/Problem_18}.

\bibitem{Don}
M.~Donahoe, \emph{Python program for A363609}, \url{https://github.com/mdonahoe/oeis/blob/main/A363609/dice.py}.

\bibitem{FC}
M.~Firsching, P.~Lezeau, S.~Mercuri, M.~Z. Horv\'ath, Y.~Dillies, C.~S\"onne, E.~Wieser, F.~Zhang, T.~Hubert, B.~Ag\"uera y Arcas, and P.~Kohli, \emph{Formal Conjectures: An Open and Evolving Benchmark for Verified Discovery in Mathematics}, arXiv:2605.13171, 2026. Repository: \url{https://github.com/google-deepmind/formal-conjectures}.

\bibitem{LW}
L.~H. Loomis and H.~Whitney, \emph{An inequality related to the isoperimetric inequality}, Bull. Amer. Math. Soc. \textbf{55} (1949), 961--962.

\bibitem{Lean4}
L.~de Moura and S.~Ullrich, \emph{The Lean 4 theorem prover and programming language}, in: Automated Deduction -- CADE 28, Lecture Notes in Comput. Sci. \textbf{12699}, Springer, 2021, 625--635.

\bibitem{Mathlib}
The mathlib Community, \emph{The Lean mathematical library}, in: Proceedings of the 9th ACM SIGPLAN International Conference on Certified Programs and Proofs (CPP 2020), ACM, 2020, 367--381.

\bibitem{A001931}
The OEIS Foundation Inc., \url{https://oeis.org/A001931} (number of fixed 3-dimensional polycubes with $n$ cells).

\bibitem{A193416}
The OEIS Foundation Inc., \url{https://oeis.org/A193416} (minimum surface area of polycubes with volume $n$), entry by J.~Barajas Martin (2011).

\bibitem{A363609}
The OEIS Foundation Inc., \url{https://oeis.org/A363609}, entry by M.~Donahoe (2023).

\end{thebibliography}

\end{document}
